% Introduction: the US rules on human authorship and AI material, the original-sin hypothesis about
% lay judgment, the rival accounts and what each implies, and what the five studies did.
% Written from scratch on 20 October 2026 and revised the same day after three reviews (law,
% fidelity, clarity); notes/results_section_checks.md, section `Introduction drafted', records the
% changes and the open questions. American spelling. Each legal source was read in an official copy
% on 20 October 2026; notes/sources.md records how. Pinpoints (also in the reference notes):
%   usco2023guidance: 88 Fed. Reg. at 16,192 (AI-determined material not human authorship, must be
%     disclaimed; prompts), 16,192-93 & nn.33-35 (selection, arrangement, modification; only the
%     human-authored aspects protected; n.35 cites 17 U.S.C. 103(b)), 16,193 (duty to disclose).
%   usco2023theatre: decision at 1 (more than de minimis AI content "must therefore be disclaimed").
%   usco2025copyrightability: Part 2 at iii and 18 (prompts alone, "given current generally
%     available technology"), 18 n.101 (the author is the party who actually creates the work),
%     24-25 & n.134 (Zarya's arrangement protectable as a compilation; edited AI output not
%     technically a derivative work, so the derivative-work rule is an analogy).
%   thaler2025perlmutter: slip op. at 2 (machine named as sole author), 3 and 9 (holding).
%   feist1991rural: 499 U.S. at 345 (originality), 364 (the quotation).
%   usco2023zarya: letter at 1 and 12 (text and selection, coordination and arrangement registered;
%     AI artwork excluded); at 10-12 no claimed image edit was registered.
%   usc2024compilations, usc2024definitions: 17 U.S.C. 103(b); 101 ("compilation", "derivative work").
%   eleventhcircuit2025pattern: Instruction 9.2 (rev. Dec. 2025) puts originality to a jury through
%     two special interrogatories. Do not widen this to "a question of fact" or a circuit split.
% Every sentence about prior work follows notes/sources.md, including its section on the novelty
% search of 20 October 2026. Do not reword the law as resting on lay psychology, Zarya as registering
% a modification, or Thaler as deciding how much AI may contribute to a human author's work. Scoped
% credit was stated after the results (argument of 14 October 2026); keep it marked as post hoc.

The US Copyright Office, which examines and records copyright claims, can register a person's contribution to a work made with artificial intelligence (AI), such as the person's own text, a creative selection and arrangement of AI material, or modifications original enough to be protected.\cite{usco2023guidance,usco2025copyrightability} Originality means independent creation with at least a minimal degree of creativity.\cite{feist1991rural} AI material beyond a trivial amount must be disclaimed, that is, excluded from the claim.\cite{usco2023guidance,usco2023theatre} The comic book \emph{Zarya of the Dawn} is registered for its text and its selection and arrangement of text and images, not its AI images.\cite{usco2023zarya} The Office takes this limit from section 103(b) of the Copyright Act: copyright in a compilation, such as \emph{Zarya}'s arrangement, or in a derivative work, one based on an earlier work, covers only what its author, the person who created it, contributed.\cite{usco2023guidance,usc2024definitions,usc2024compilations} It treats edited AI images as derivative works only by analogy, because in its view wholly AI-made images are not works of authorship.\cite{usco2025copyrightability} Copyright requires a human author: in \emph{Thaler v.\ Perlmutter}, which concerned a machine named as sole author, a federal appeals court held that the Act requires one.\cite{thaler2025perlmutter,usco2025copyrightability} In 2025 the Office concluded that, with the technology then available, prompts alone gave a user too little control to author the output.\cite{usco2023guidance,usco2025copyrightability} These rules rest on no claim about how people judge AI-made work, yet they ask whoever applies them to judge a person's part on its own while knowing the base was machine-made: examiners, who read applicants' disclosures of AI content,\cite{usco2023guidance} and jurors, because originality can go to a jury in an infringement suit.\cite{eleventhcircuit2025pattern}

We call original sin the hypothesis that lay judgment does not separate the two parts: that an AI origin acts like a stain passing to whatever is built from it. It concerns people, not law; our debriefing wrongly told participants that a categorical version was law (Methods). It extends to AI the belief, studied as contagion, that a person's essence passes into objects they touched or merely designed.\cite{newman2011celebrity,stavrova2016contamination} The same image is rated lower when labeled AI-made than human-made,\cite{bellaiche2023humans,horton2023bias,millet2023defending} a drop we call the label's penalty, small to moderate in a meta-analysis\cite{derooij2025bias} and not found in every design.\cite{neef2025drawing,messingschlager2025mind} Work described as made by a person with AI is also rated lower on most measures, even a painting made by hand from AI-supplied compositions, though its artist was rated more creative.\cite{horton2023bias,messer2024cocreating} These studies rated the finished work or its maker. Apart from our earlier report of two of the studies below,\cite{kamper2025originalsin} we know of no study that showed an image beside a person's edit of it, varied only its described origin, and asked about the resulting second work.

On the effort-signal account, the label lowers ratings of originals only because it signals little work, as AI-labeled work is judged to take less effort and time;\cite{bellaiche2023humans,horton2023bias,magni2024humans} a stated 100 hours should then shrink the penalty to almost nothing. If the label instead says the act was prompting, not drawing, the penalty should persist. Because copyright rewards ``originality, not effort,''\cite{feist1991rural} we ask whether the label lowers effort ratings more than creativity ratings. For second works, Study~2's registration (a dated record of its hypotheses and analyses, not a copyright registration) sets out three accounts, which we call taint, dilution and severance. Taint is original sin applied to second works: the same penalty at every level of modification, at full strength as large as on originals. Under dilution, raters judge the second work partly by its original, so the penalty shrinks as the edit grows. Under severance, raters judge only what the person added, so there is no penalty, which parallels section 103(b). After seeing the results, we stated a principle, scoped credit: people's ratings credit a person for what they take that person to have done, and a label says which act produced what they see. It implies a penalty on originals, whose making the label describes; little carry-over to second works, whose making both groups read about alike; and an effect of stated hours only when they describe the rated act.

A label told each participant either that the artist, Alex, had made the originals with design and sketching tools or that he had prompted an AI system to make them; all were in fact made with OpenAI's DALL\textperiodcentered E through ChatGPT, so any group difference is a label effect. Studies 1 and 3 asked about the originals, Study~3 stating a time of under 1, 10 or 100 hours. In Studies 2, 4 and 5, both groups read, in the same words, that Alex made three second works from each original with his design and sketching tools: None, the original unchanged; Slight, a light edit; and Dramatic, a heavy one. We made the edits by hand. Study~4 also stated hours for each original, and Study~5 hours for each edit. Participants saw each second work beside its original before rating it. The questions concerned the new work and the edit, not how the original was made, so any label effect is a carry-over; its absence would parallel section 103(b) but would also follow if raters answered the question asked (Discussion). All five studies were registered on 8 October 2026, after data collection, so none was preregistered; only Study~5 was registered before any of its ratings had been averaged or compared by label. The numbering follows the argument; Studies 2 and 4 ran first (Methods).
